\begin{figure}[t]
\centering
\figplaceholder[2.3in]{\columnwidth}{图 3——作用于运动目标（示意图）}{%
\textbf{版式：}单行三个等宽子图 (a)--(c)，矢量示意图，使用论文的语义配色；各子图使用同一机械臂与 T 形零件；不使用照片。
\textbf{(a) 同一位姿，不同运动：}零件的两段历史 A（\textcolor{obscol}{蓝色实线}）与 B（\textcolor{obscol}{蓝色虚线}），在 $-100$、$-50$、$0$\,ms 处画出半透明位姿，二者在采集时刻 $\tau_t$ 到达同一位姿。标注：“仅凭当前位姿无法区分”。
\textbf{(b) 执行时刻的位姿：}子图下方一条时间轴：从采集 $\tau_t$ 起每隔 $\Delta_a=50$\,ms 一个刻度，对应动作块各指令的时刻 $\tapp_{t,j}$；动作块在 $\tau_t+\ell_t$ 到达之前的刻度灰显为已过期，第一条被执行的指令标为 $\tapp$。上方画出 A 与 B 在 $\tapp$ 时位置和朝向都不同的零件（\textcolor{predcol}{橙色}三轴）。标注：“指令必须适合 $\tapp$ 时的位姿”。
\textbf{(c) 角色：}同一场景中的三类物体。\emph{目标}：零件，标出末端执行器坐标系 $E$、目标坐标系 $O$、接触点与容差 $\varepsilon$（mm）。\emph{环境中的障碍物}：容器壁或遮挡物，以箭头标出间隙 $m\gg\varepsilon$。\emph{环境中的驱动物}：承载零件的传送带及其运动箭头，可另画一个正在接近的撞击器；标注“推动目标运动”。\emph{背景}：桌沿与远处物体，以\textcolor{gtcol}{灰色}去饱和显示；标注“不影响任务”。三行图例：目标 $\to$ 精确 SE(3)；环境 $\to$ 粗粒度 3D 流；背景 $\to$ 无需建模。
\textbf{规则：}概念示意，位置为假设；\textcolor{obscol}{蓝色}为过去，\textcolor{predcol}{橙色}为未来，\textcolor{actcol}{紫色}为指令，\textcolor{gtcol}{灰色}为背景。}
\caption{\textbf{作用于运动目标}（示意图）。(a) 两段历史在采集时刻 $\tau_t$ 到达相同位姿，但对应不同运动。(b) 在执行时刻 $\tapp$，它们对应不同的目标位姿。(c) 接触要求末端执行器坐标系 $E$ 与目标 $O$ 在 $\varepsilon$ 内对齐；环境中的障碍物只通过它留下的间隙 $m\gg\varepsilon$ 起作用，传送带这类驱动物通过传给目标的运动起作用，背景则不影响任务。}
\label{fig:principles}
\end{figure}